\chapter{Time-Lapse Detection of Vertical Movement}
\label{ch:tl-vertical}

This chapter mirrors \cref{ch:tl-horizontal}, replacing lateral displacement
of the scatterer with \emph{vertical} (depth) displacement. The two
directions are not expected to behave identically: \cref{sec:th-duality}
showed that a migrated image separates lateral spatial frequencies like a
prism but does not separate vertical ones, so the amplitude-based
detectability established here is an important point of comparison for the
phase-based vertical results of \cref{sec:tlp-vertical}.

\section{Setup: Moving Scatterer Scenarios}

A single PEC cylinder ($r=\SI{28}{mm}$) sits at lateral position
$x=\SI{2.0}{m}$ and is displaced \emph{downward} from its baseline depth
across seven scenarios, from $1\lambda$ down to $\tfrac{1}{32}\lambda$
(\cref{fig:vtl-setup}), using the same domain and grid as
\cref{ch:resolution,ch:tl-horizontal}.

\begin{figure}[htbp]
  \centering
  \begin{subfigure}[t]{0.48\textwidth}
    \includegraphics[width=\linewidth]{VTL_001_Vertical_TimeLapse_Study__Model_Geometry__domain_4010_m_Δx.png}
    \caption{Model geometry}
  \end{subfigure}
  \hfill
  \begin{subfigure}[t]{0.48\textwidth}
    \includegraphics[width=\linewidth]{VTL_002_Moving_scatterer__all_7_scenarios__r__28_mm_x__20_mBaseline.png}
    \caption{All 7 vertical displacement scenarios}
  \end{subfigure}
  \caption{Forward-model setup for the vertical time-lapse study: (a) the
    gprMax domain and grid; (b) the vertical displacement of the scatterer
    across all seven scenarios, $r=\SI{28}{mm}$, fixed at $x=\SI{2.0}{m}$.}
  \label{fig:vtl-setup}
\end{figure}

\section{Raw B-Scan Generation}

\Cref{fig:vtl-bscans} shows the background-baseline and vertically
time-lapsed raw B-scans and the background-subtracted result, with the
expected arrival times for each scenario marked.

\begin{figure}[htbp]
  \centering
  \begin{subfigure}[t]{0.48\textwidth}
    \includegraphics[width=\linewidth]{VTL_003_GPR_B-Scans__Background_Baseline_and_Vertical_TimeLapsed_Mod.png}
    \caption{Background \& time-lapsed}
  \end{subfigure}
  \hfill
  \begin{subfigure}[t]{0.48\textwidth}
    \includegraphics[width=\linewidth]{VTL_004_GPR_B-Scans__Background_Subtracted__green_dashed__expected_a.png}
    \caption{Background-subtracted}
  \end{subfigure}
  \caption{Raw B-scans for the vertical time-lapse study: (a) background and
    time-lapsed models; (b) background-subtracted, with expected arrival
    times marked (green, dashed).}
  \label{fig:vtl-bscans}
\end{figure}

\section{Signal Conditioning}

\Cref{fig:vtl-taper} shows the effect of the standard tapering and $t_0$-shift
conditioning (\cref{sec:meth-conditioning}) on the $1\lambda$ dataset.

\begin{figure}[htbp]
  \centering
  \begin{subfigure}[t]{0.48\textwidth}
    \includegraphics[width=\linewidth]{VTL_005_Effect_of_Tapering_and_t0_Shift__1λ_dataset_single_trace.png}
    \caption{Single trace}
  \end{subfigure}
  \hfill
  \begin{subfigure}[t]{0.48\textwidth}
    \includegraphics[width=\linewidth]{VTL_006_B-scan_effect_of_tapering_and_t0_shift__1λ_dataset.png}
    \caption{Full B-scan}
  \end{subfigure}
  \caption{Effect of tapering and the $t_0$ shift on the $1\lambda$ vertical
    displacement dataset: (a) a single trace; (b) the full B-scan.}
  \label{fig:vtl-taper}
\end{figure}

\section{Kirchhoff Migration and Time-Lapse Differencing}

\Cref{fig:vtl-kirchhoff} shows the Kirchhoff-migrated image for all seven
vertical displacement scenarios and the resulting time-lapse difference.

\begin{figure}[htbp]
  \centering
  \begin{subfigure}[t]{0.48\textwidth}
    \includegraphics[width=\linewidth]{VTL_007_Kirchhoff_Migration__All_7_Datasets____f_c15_GHz____aperture.png}
    \caption{All 7 scenarios}
  \end{subfigure}
  \hfill
  \begin{subfigure}[t]{0.48\textwidth}
    \includegraphics[width=\linewidth]{VTL_008_Kirchhoff_Migration_zoomed____f_c15_GHz____aperture40.png}
    \caption{Zoomed}
  \end{subfigure}

  \vspace{0.6em}
  \begin{subfigure}[t]{0.48\textwidth}
    \includegraphics[width=\linewidth]{VTL_009_Kirchhoff_Migration__TimeLapse_Differences_migrated__migrate.png}
    \caption{Time-lapse difference}
  \end{subfigure}
  \hfill
  \begin{subfigure}[t]{0.48\textwidth}
    \includegraphics[width=\linewidth]{VTL_010_Kirchhoff_Migration__TimeLapse_Differences_zoomed.png}
    \caption{Difference, zoomed}
  \end{subfigure}
  \caption{Kirchhoff migration of the vertical time-lapse study
    ($f_c=\SI{15}{\GHz}$, aperture $=40$): (a,b) migrated image for all
    scenarios and zoomed; (c,d) time-lapse difference and zoomed.}
  \label{fig:vtl-kirchhoff}
\end{figure}

\section{Gazdag Migration and Time-Lapse Differencing}

\Cref{fig:vtl-gazdag} repeats the analysis with Gazdag phase-shift migration.

\begin{figure}[htbp]
  \centering
  \begin{subfigure}[t]{0.48\textwidth}
    \includegraphics[width=\linewidth]{VTL_011_Gazdag_Phase-Shift_Migration__All_7_Datasets____f_c15_GHz.png}
    \caption{All scenarios}
  \end{subfigure}
  \hfill
  \begin{subfigure}[t]{0.48\textwidth}
    \includegraphics[width=\linewidth]{VTL_012_Gazdag_Phase-Shift_Migration_zoomed____f_cf_c_GHz_GHz.png}
    \caption{Zoomed}
  \end{subfigure}

  \vspace{0.6em}
  \begin{subfigure}[t]{0.48\textwidth}
    \includegraphics[width=\linewidth]{VTL_013_Gazdag_Migration__TimeLapse_Differences_migrated__migrated_b.png}
    \caption{Time-lapse difference}
  \end{subfigure}
  \hfill
  \begin{subfigure}[t]{0.48\textwidth}
    \includegraphics[width=\linewidth]{VTL_014_Gazdag_Migration__TimeLapse_Differences_zoomed.png}
    \caption{Difference, zoomed}
  \end{subfigure}
  \caption{Gazdag phase-shift migration of the vertical time-lapse study
    ($f_c=\SI{15}{\GHz}$): (a,b) migrated image and zoomed; (c,d) time-lapse
    difference and zoomed.}
  \label{fig:vtl-gazdag}
\end{figure}

\section{Back-Propagation Migration and Time-Lapse Differencing}

\Cref{fig:vtl-backprop} shows the back-propagation result for both the
field-magnitude and $E_z$ images, focused at $t=\SI{1906}{\ns}$.

\begin{figure}[htbp]
  \centering
  \begin{subfigure}[t]{0.48\textwidth}
    \includegraphics[width=\linewidth]{VTL_015_Back-Propagation_E__All_7_Datasets____focus_at_1906_ns.png}
    \caption{$\lVert\mathbf{E}\rVert$, all scenarios}
  \end{subfigure}
  \hfill
  \begin{subfigure}[t]{0.48\textwidth}
    \includegraphics[width=\linewidth]{VTL_016_Back-Propagation_E_zoomed____focus_at_1906_ns.png}
    \caption{$\lVert\mathbf{E}\rVert$, zoomed}
  \end{subfigure}

  \vspace{0.6em}
  \begin{subfigure}[t]{0.48\textwidth}
    \includegraphics[width=\linewidth]{VTL_017_Back-Propagation_Ez__All_7_Datasets____focus_at_1906_ns.png}
    \caption{$E_z$, all scenarios}
  \end{subfigure}
  \hfill
  \begin{subfigure}[t]{0.48\textwidth}
    \includegraphics[width=\linewidth]{VTL_018_Back-Propagation_Ez_zoomed____focus_at_1906_ns.png}
    \caption{$E_z$, zoomed}
  \end{subfigure}

  \vspace{0.6em}
  \begin{subfigure}[t]{0.48\textwidth}
    \includegraphics[width=\linewidth]{VTL_019_Back-Propagation__TimeLapse_Differences_Ez_Ez__Ez_baseline.png}
    \caption{$E_z$ time-lapse difference}
  \end{subfigure}
  \hfill
  \begin{subfigure}[t]{0.48\textwidth}
    \includegraphics[width=\linewidth]{VTL_020_Back-Propagation__TimeLapse_Differences_Ez_zoomed.png}
    \caption{Difference, zoomed}
  \end{subfigure}
  \caption{Time-reversal back-propagation migration of the vertical
    time-lapse study, focused at $t=\SI{1906}{\ns}$: field-magnitude (a,b)
    and $E_z$ (c,d) images, and the $E_z$ time-lapse difference (e,f).}
  \label{fig:vtl-backprop}
\end{figure}

\section{Vertical Detectability Summary}
\label{sec:vtl-detectability}

\Cref{fig:vtl-summary} overlays the signed time-lapse-difference amplitude
from all three algorithms and plots the normalised \emph{vertical}
point-spread function of that difference at $x=\SI{2.0}{m}$, as a function
of vertical displacement.

\begin{figure}[htbp]
  \centering
  \begin{subfigure}[t]{0.48\textwidth}
    \includegraphics[width=\linewidth]{VTL_021_Vertical_TimeLapse_Migration_Comparison__Signed_Amplitude.png}
    \caption{Signed-amplitude comparison}
  \end{subfigure}
  \hfill
  \begin{subfigure}[t]{0.48\textwidth}
    \includegraphics[width=\linewidth]{VTL_022_Normalised_Vertical_PSF__TimeLapse_Difference_at_x__20_m.png}
    \caption{Normalised vertical PSF of the difference}
  \end{subfigure}
  \caption{(a) Signed time-lapse-difference amplitude for all three
    migration algorithms; (b) the normalised vertical PSF of the difference
    image at $x=\SI{2.0}{m}$, swept across vertical displacements from
    $1\lambda$ to $\tfrac{1}{32}\lambda$.}
  \label{fig:vtl-summary}
\end{figure}

\draftnote{compare the smallest reliably-detected vertical displacement in
\cref{fig:vtl-summary}b against the lateral result of
\cref{fig:tl-summary}b (\cref{sec:tl-detectability}) and discuss whether the
lateral/vertical asymmetry predicted by the instantaneous-phase theory of
\cref{sec:th-duality} is also visible at the amplitude level, or only
emerges once the phase-plane estimator of \cref{sec:tlp-vertical} is
applied.}
